\documentclass[11pt]{article}
\usepackage[margin=1in]{geometry}
\usepackage{booktabs}
\usepackage{longtable}
\usepackage{hyperref}
\usepackage{xcolor}
\usepackage{enumitem}
\usepackage{amsmath}
\usepackage{tabularx}

\title{Replication Report: Delgado-Su\'arez et al.\ (2021)\\
\large Genomic surveillance of antimicrobial resistance shows cattle and poultry\\
are a moderate source of multi-drug resistant non-typhoidal \emph{Salmonella} in Mexico}
\author{Ollie (OpenClaw AI) --- BV-BRC Replication Project, target BVBRC-31}
\date{2026-07-01}

\begin{document}
\maketitle

\begin{center}
\textbf{Verdict: PARTIAL REPLICATION}
\end{center}

\noindent\textbf{Paper:} Delgado-Su\'arez EJ, et al. \emph{PLoS ONE} 16(5):e0243681 (2021).\\
\textbf{DOI:} \href{https://doi.org/10.1371/journal.pone.0243681}{10.1371/journal.pone.0243681} \quad
\textbf{PMID:} 33951039 \quad \textbf{PMCID:} PMC8099073 \quad \textbf{Open access:} yes (CC0)\\
\textbf{Data:} BioProject \textbf{PRJNA480281} (77 study isolates); per-isolate accessions in S1 File.

\section{Paper}

The study phenotypically and genotypically characterizes AMR in \textbf{77 non-typhoidal \emph{Salmonella} (NTS)}
isolates from Mexican \textbf{bovine lymph nodes (n=48)} and \textbf{ground beef (n=29)}, and compares their
genotypic AMR to \textbf{2{,}400 public Mexican NTS genomes}. Genotypic pipeline (Methods): assemble reads
$\to$ \textbf{SeqSero v1.2} serovar prediction $\to$ \textbf{MLST} $\to$ \textbf{AMRFinderPlus v3.8.4}.
Headline claims: tetracycline is the most common resistance (40.3\%), \textbf{26\% are MDR}, MDR is more likely
in ground beef than lymph nodes ($\chi^2=12.0$, $P=0.0005$), MDR is serovar-associated ($\chi^2=24.5$, $P<0.0001$)
with \textbf{Typhimurium = 40\% of MDR strains}, and \textbf{9/10 Typhimurium carry SGI1} with a class-1 integron
(aadA2, blaCARB-2, floR, sul1, tetG = penta-resistant). ramR mutations were associated with MDR
($\chi^2=17.7$, $P<0.0001$).

\section{Claims tested}

\begin{longtable}{clcc}
\toprule
\# & Claim & Type & Tested \\
\midrule
C1 & 77 isolates (48 LN + 29 GB) publicly available via PRJNA480281 & Data availability & Yes \\
C2 & Serovar distribution reproducible by in-silico typing & Genomic & Yes \\
C3 & Per-antibiotic-class AMR prevalence (tet top; ceph/carbapenem rare) & Genomic & Yes \\
C4 & $\sim$26\% of isolates are MDR ($\geq$3 classes) & Genomic & Yes \\
C5 & MDR higher in ground beef vs lymph nodes ($\chi^2=12.0$, $P=0.0005$) & Statistical & Direction / partial \\
C6 & Typhimurium = 40\% of MDR strains & Genomic & Yes \\
C7 & 9/10 Typhimurium carry SGI1 penta-resistance cassette & Genomic & Yes \\
C8 & 100\% QRDR mutations; ramR--MDR association & Genomic (SNP) & Attempted; not reproduced \\
\bottomrule
\end{longtable}

\section{Method}

Heavy compute on \textbf{uicgpu} (8$\times$A100, 255 cores, 2\,TB RAM). All tools free/open; LLM judge via
free Argo proxy.

\begin{enumerate}[leftmargin=*]
  \item \textbf{Paper + accessions.} Europe PMC full text named BioProject PRJNA480281; S1 File (\texttt{s001.xlsx})
        provided the 77 isolates. Parsed cohort matches paper: 48 LN + 29 GB; serovars Anatum 23, Reading 22,
        Typhimurium 10, London 9, Kentucky 6, Fresno 4, others.
  \item \textbf{BioSample $\to$ assembly.} NCBI Datasets v2alpha for all 77: 68 had GenBank (GCA) assemblies;
        9 (a newer Reading batch + 2 others) left for a future reads-based pass.
  \item \textbf{Genome download.} \texttt{datasets download genome accession -{}-inputfile assembly\_list.txt -{}-include genome}
        $\to$ 68 FASTAs.
  \item \textbf{Genotyping} (bioconda: AMRFinderPlus 3.12.8 / DB 2024-07-22.1, SeqSero2 1.3.2, mlst 2.35.0):
    \begin{itemize}
      \item \texttt{amrfinder -n <fna> -{}-organism Salmonella -{}-plus} (per isolate, 16-way parallel).
      \item \texttt{SeqSero2\_package.py -m k -t 4 -i <fna>}.
      \item \texttt{mlst assemblies/*.fna}.
    \end{itemize}
  \item \textbf{Analysis} (\texttt{work/analyze.py}): AMR gene $\to$ class from AMRFinder's \texttt{Class} column
        (core intrinsic efflux mdsAB/golST excluded); MDR = $\geq$3 acquired classes (Magiorakos 2012);
        per-class prevalence; MDR by source with 2$\times$2 $\chi^2$; Typhimurium SGI1 penta-set membership;
        serovar concordance vs paper.
  \item \textbf{Verdict} (\texttt{work/judge.py}): LLM judge = Argo \texttt{argo:gpt-5.2} (free).
\end{enumerate}

\section{Results vs Paper}

\paragraph{C1 --- Cohort availability.}
S1 File yields exactly 77 isolates (48 LN + 29 GB); 68 have GenBank assemblies (88\%). Serovar counts from
S1 match paper text one-for-one.

\paragraph{C2 --- Serovar typing (SeqSero2).}
\textbf{67/68 (98.5\%)} concordant; the single nominal miss (GCA\_007741155.1, ``Reading'') is SeqSero2
reporting the Reading antigenic formula \texttt{I -:e,h:1,5}, which \emph{is} Reading, so effective
concordance is 68/68.

\paragraph{C3 --- Per-class AMR prevalence (68 isolates).}
\begin{center}
\begin{tabular}{lccl}
\toprule
Class & Paper (pheno \%) & This work (geno \%) & Note \\
\midrule
Tetracycline           & 40.3 & 33.8 (23/68) & Top acquired class --- matches paper ranking \\
$\beta$-lactam         & 26.0 / 20.8 & 11.8 (8/68) & Genotypic $<$ phenotypic (some intrinsic AmpC) \\
Phenicol               & 19.5 & 10.3 (7/68)  & floR-driven; SGI1 subset \\
Sulf.+trimethoprim     & 16.9 & 11.8 (8/68)  & sul1-driven \\
Aminoglycoside         & ---  & 13.2 (9/68)  & aadA2 etc. \\
Quinolone (PMQR)       & $>$55 dec.\ suscept. & 47.1 (32/68) & qnrB19 in 31/68 \\
Cephalosporin (ESBL)   & 0--9 & 0 (0/68)     & Rare/absent, matches paper \\
\bottomrule
\end{tabular}
\end{center}

\paragraph{C4 --- MDR prevalence.}
\textbf{16/68 = 23.5\% MDR} vs paper \textbf{26\%}. Difference 2.5\,pp, within expected geno-vs-pheno +
68/77 coverage variance. \textbf{Strong match.}

\paragraph{C5 --- MDR by source.}
Ground beef 8/24 = 33.3\%; lymph nodes 8/44 = 18.2\%. \textbf{Direction reproduced}, but $\chi^2=1.98$
($p=0.16$) on the 68-subset --- \emph{not} significant, whereas the paper reports $\chi^2=12.0$, $P=0.0005$.
Attributable to (a) 9 missing isolates, (b) genotypic-vs-phenotypic MDR classification.

\paragraph{C6 --- Typhimurium share of MDR.}
6/16 MDR are Typhimurium = \textbf{37.5\%} vs paper \textbf{40\%}. Of 7 Typhimurium analyzed, 6 are MDR.
MLST corroborates: 7 isolates are ST19 (canonical Typhimurium).

\paragraph{C7 --- SGI1 penta-resistance cassette.}
\textbf{6/7 Typhimurium} carry the complete aadA2 / blaCARB-2 / floR / sul1 / tetG set --- exactly the paper's
cassette. Paper reports 9/10; only 7 Typhimurium had assemblies here, so 6/7 (86\%) is the direct analogue
of 9/10 (90\%). Clean direct replication of the central resistance-mechanism claim.

\paragraph{C8 --- Point mutations (QRDR / ramR).}
AMRFinderPlus \texttt{-{}-organism Salmonella} (mutation search confirmed in logs) emitted \textbf{zero}
curated resistance point mutations across all 68 isolates. The paper's 100\% QRDR (gyrA/gyrB/parE), soxRS,
and ramR--MDR results reflect \emph{raw sequence comparison vs a reference}, not AMRFinder's curated
resistance-SNP catalog. Under a standard curated catalog this claim does not reproduce --- a tool-methodology
difference, not a data problem.

\paragraph{MLST corroboration.}
ST distribution: ST64 $\times$21 (Anatum), ST1628 $\times$19 (Reading), ST155 $\times$8 (London),
ST19 $\times$7 (Typhimurium), ST649 $\times$4, ST198 $\times$4 (Kentucky), others. Serovar$\leftrightarrow$ST
pairing is internally consistent.

\section{Verdict}

\textbf{PARTIAL REPLICATION.}\\
\emph{Reproduced:} cohort (77 = 48 LN + 29 GB); serovar distribution (67/68 $\to$ effectively 68/68);
MDR 23.5\% vs 26\%; tetracycline-topped class ranking; ceph/carbapenem absent; Typhimurium 37.5\% of MDR;
\textbf{SGI1 penta-cassette in 6/7 Typhimurium.}\\
\emph{Not reproduced:} statistical significance of GB-vs-LN MDR difference (direction correct, $p$ n.s.\ on
68-subset); QRDR/ramR point-mutation claims (curated-catalog vs custom-SNP-calling).\\
\emph{LLM judge (Argo gpt-5.2):} verdict \textbf{PARTIAL}, coverage 7/10, agreement 5/10.

\section{GENUINE CRITIQUE}

This section is deliberately unsanitized. It is what an honest replicator would flag to a program officer.

\subsection*{What clearly holds up}
\begin{itemize}[leftmargin=*]
  \item \textbf{Cohort and data availability are exemplary.} 77/77 isolates traceable through S1 File to
        BioSample/SRR; 68/77 (88\%) have curated GenBank assemblies. This is above the median for
        veterinary-genomics replication targets in the BV-BRC set.
  \item \textbf{Serovar calls are essentially perfect} (68/68 after resolving the antigenic-formula naming).
        The genotypic pipeline the paper describes is truly deterministic and portable.
  \item \textbf{The central mechanistic claim reproduces cleanly.} 6/7 Typhimurium carry the exact
        aadA2/blaCARB-2/floR/sul1/tetG SGI1 penta-cassette. This is not a hand-wave --- it is the same five
        genes in the same isolates, called by a different investigator with a different tool version.
        This is the strongest single evidence in the study and it survives independent replication.
  \item \textbf{MDR prevalence sits within 2.5\,pp} (23.5 vs 26). For a genotype-vs-phenotype comparison
        that is arguably as close as one can honestly expect.
\end{itemize}

\subsection*{Where the paper is doing more than the sequence data alone supports}
\begin{itemize}[leftmargin=*]
  \item \textbf{The GB-vs-LN significance ($\chi^2=12.0$, $P=0.0005$) is fragile.} On the 68-genome subset
        the direction survives but the significance evaporates ($p=0.16$). Nine missing isolates should not
        collapse a $P=0.0005$ result if the effect were as robust as claimed --- this suggests the paper's
        significance is at least partly power-driven, and possibly leans on the AST-based (phenotypic)
        MDR call being more lenient than a genotype-based one.
  \item \textbf{The QRDR/ramR mutation claims do not reproduce under standard curated catalogs.} The paper
        reports \emph{100\%} QRDR mutations and a $\chi^2=17.7$ ramR--MDR association. AMRFinderPlus with
        the paper's own organism flag returns \emph{zero} curated resistance point mutations across all 68
        isolates. The paper is using a raw SNP-vs-reference approach (call every non-synonymous change) and
        then labelling it ``resistance mutation''. That is a legitimate methodological choice but it is not
        the same construct as ``curated, resistance-linked SNP'' and readers should not conflate the two.
        The paper does not make this distinction cleanly.
  \item \textbf{Genotype--phenotype gap is honestly present.} $\beta$-lactam is 26\% phenotypic vs
        11.8\% genotypic; phenicol 19.5 vs 10.3; SXT 16.9 vs 11.8. The rankings match, but the paper's
        rhetoric ($\sim$``genomic prediction of resistance'') implicitly assumes tight geno$\to$pheno
        concordance, which the numbers themselves do not support at the per-class level.
  \item \textbf{The 2{,}400-genome public comparison was not testable here} (S2 File accessions are
        available but this replication did not run them). The paper's headline framing ---
        ``cattle and poultry are a moderate source of MDR NTS'' --- rests on the 2{,}400-set contrast, and
        that specific sentence is therefore \emph{not} independently checked in this report.
\end{itemize}

\subsection*{What the paper does not do that would strengthen it}
\begin{itemize}[leftmargin=*]
  \item \textbf{No confidence intervals} on the 26\% MDR estimate; the CI on 20/77 is roughly $\pm$10\,pp,
        which would deflate the ``moderate source'' framing.
  \item \textbf{No plasmid backbone typing} beyond the SGI1 island. If the class-1 integron is on IncA/C or
        IncHI2 that is a very different public-health story (mobilizable to \emph{E. coli}) than if it is
        chromosomal.
  \item \textbf{No temporal analysis.} All 77 are collected in a narrow window; the 2{,}400-set comparison
        conflates decades of surveillance with a single 2018--2019 slice, which is not corrected for in the
        paper's $\chi^2$ tests.
  \item \textbf{Reproducibility footprint is above average for the field but still not one-click.} S1 File
        is XLSX, not FASTA-linked; the mutation call is not scripted; there is no container image or Snakemake
        pipeline archived alongside the paper.
\end{itemize}

\subsection*{Bottom line}
The paper's \emph{genotypic core} (serovar distribution, MDR magnitude, Typhimurium's disproportionate MDR
share, SGI1 penta-cassette) is real and independently reproducible. The paper's \emph{statistical and
mutational claims} (GB-vs-LN significance, 100\% QRDR, ramR--MDR association) are weaker under independent
replication: the significance is power-fragile and the mutation calls are catalog-dependent in a way the paper
does not disclose. A conscientious reader should trust the mechanism (SGI1 in Typhimurium) more than the
inference (cattle/beef as a distinct source), and should read ``100\% QRDR mutations'' as
``100\% carry a non-synonymous QRDR SNP,'' not ``100\% predicted quinolone resistance.''

\section{Coverage / Agreement}

\begin{itemize}[leftmargin=*]
  \item \textbf{Coverage: 7/10.} C1, C2, C3, C4, C6, C7 fully tested and reproduced; C5 direction tested;
        C8 tested but not reproduced. Not attempted: the 2{,}400-genome public comparison and the 9 missing
        isolates.
  \item \textbf{Agreement: high on what was end-to-end testable.} Serovar 67/68, MDR 23.5 vs 26,
        Typhimurium 37.5 vs 40, SGI1 6/7 vs 9/10. Disagreements confined to source-attribution significance
        (power) and point-mutation calling (catalog).
\end{itemize}

\section{Path to full REPLICATED}

\begin{enumerate}[leftmargin=*]
  \item Assemble the 9 missing isolates de novo from SRR reads (SKESA/SPAdes) $\to$ recover 9/10 SGI1 count
        and restore statistical power for C5.
  \item Reproduce paper's raw-SNP approach on gyrA/gyrB/parC/parE/ramR/soxRS aligned to \emph{S.} Typhimurium
        LT2 --- not just the curated resistance-SNP catalog.
  \item Pull the 2{,}400 S2-File genomes and rerun AMRFinder to test the ``cattle/poultry highest MDR'' claim.
  \item Obtain the phenotypic AST table for a like-for-like geno-vs-pheno comparison.
\end{enumerate}
None require paid resources, GPUs, or restricted data --- only additional CPU time.

\end{document}
